We can now define the adaptivity of a program formally. This notion is formulated in terms of an initial trace, specifying the value of the input variables, as the walk on the graph $\traceG({c})$, which has the largest number of query requests.


\begin{defn}[Walk]
\label{def:finitewalk}
Given a well-formed program $c$ with its semantics-based dependency graph $\traceG({c}) = (\traceV, \traceE, \traceW, \traceF)$ of a program $c$, a \emph{walk} $k$ is a function that maps an initial trace $\trace_0$ to a sequence of vertices $(v_1, \ldots, v_{n})$
for which there is a sequence of edges $(e_1 \ldots e_{n - 1})$  satisfying
\begin{itemize}
\item $e_i = (v_{i},v_{i + 1}) \in \traceE$ for every $1 \leq i < n$,
\item and $v_i$ appears in $(v_1, \ldots, v_{n})$ at most $w_i(\trace_0)$ times for every $v_i \in \traceV$ and $(v_i, w_i) \in \traceW$.  
\end{itemize}
$k(\trace_0) = (v_1, \ldots, v_{n})$
%  and the length of $k(\trace_0)$ is the number of vertices in its vertex sequence, i.e., $|k(\trace_0)| = n$.}
% The length of $k(\trace_0)$ is the number of vertices in its vertex sequence, i.e., $\len(k)(\trace_0) = n$.
and we denote by $\walks(\traceG(c))$
the set of all the walks $k$ in $\traceG(c)$.
\end{defn} 
Because for the adaptivity
% is intuitively 
we are interested in the dependency between queries,
we calculate a special ``length'' of a walk, the \emph{query length},  by counting only the vertices
corresponding to queries.
\begin{defn}[Query Length]
\label{def:qlen}
Given 
the semantics-based dependency graph $\traceG({c})$ of a well-formed program $c$,
 and a \emph{walk} 
 $k \in \walks(\traceG(c))$ 
%  with vertex sequence $(v_1, \ldots, v_{n})$, 
the \emph{query length} of $k$ is a function $\qlen(k):\ftdom_0(c) \to \mathbb{N}$ that 
% given an initial trace $\trace_0$ 
gives
the number of vertices that correspond to query variables in the vertex sequence of $k(\trace_0)$
for any initial trace $\trace_0 \in \ftdom_0(c)$ as follows, 
\begin{center}
   $
  \qlen(k) = \lambda \trace_0 \sthat |\big( v \mid v \in (v_1, \ldots, v_{n}) \land (v, 1) \in \traceF(c) 
%   \land
%   k \in \walks(\traceG(c)) 
  \land k(\trace_0) = (v_1, \ldots, v_{n}) \big)|,
$
\end{center}
where the notation $| (\ldots) |$ gives the number of vertices in a sequence.
\end{defn}
Then we define for a well-formed program its adaptivity as a function as below.
\begin{defn}
    [Adaptivity of a Program]
    \label{def:trace_adapt}
    Given a well-formed program ${c}$, 
    its adaptivity $A(c)$ is a function 
    $A(c) : \ftdom_0(c)\to \mathbb{N}$ such that for an
    initial trace $\trace_0 \in \ftdom_0(c)$, 
\begin{center}
$
    A(c) = \lambda \trace_0 \sthat \max \big 
    \{ \qlen(k)(\trace_0) \mid k \in \walks(\traceG(c)) \big \} 
$
\end{center}
\end{defn}
As discussed above, programs that has non-terminating behaviors do not have the intuitive adaptivity. So we define the adaptivity only for the well-formed program.